
\subsubsection*{2.4.A}
\phantomsection
\label{sga3:VIIB:2.4.A}

Let \(G\) and \(Q\) be formal \(k\)-groups that are topologically
flat. Suppose that there are homomorphisms
\[
  \sigma:Q\longrightarrow G
  \qquad\text{and}\qquad
  \pi:G\longrightarrow Q
\]
such that
\[
  \pi\circ\sigma=\operatorname{id}_Q.
\]
\footnote{\label{sga3:VIIB:nde:92}The editors added paragraphs
2.4.A and \SGARefLink{sga3:VIIB:2.4.B}{2.4.B}, which will be useful in \SGARefLink{sga3:VIIB:5.1.3}{5.1.3} and \SGARefLink{sga3:VIIB:2.9}{2.9},
respectively.}
In particular, \(\sigma\) is a monomorphism, so \(Q\) is a formal
subgroup of \(G\) \textup{(cf.\ \SGARefLink{sga3:VIIB:1.3}{Remark 1.3})}. Let
\[
  N=\operatorname{Ker}(\pi)
\]
and let \(\sigma'\) be the inclusion \(N\hookrightarrow G\).
Then \(G\) is the semidirect product of \(N\) by \(Q\)
\textup{(cf.\ I, \SGARefLink{sga3:i:proposition:2-3-8}{2.3.8})}. In other words, for every
\[
  B\in\operatorname{Ob}(\operatorname{Alf}/k),
\]
identify \(N(B)\) and \(Q(B)\) with their images in \(G(B)\) under
\(\sigma'\) and \(\sigma\). Then \(N(B)\) is a normal subgroup of
\(G(B)\), and the map
\begin{equation}
  \mu:N(B)\times Q(B)\longrightarrow G(B),
  \qquad
  (x,q)\longmapsto xq
  \tag{1}
  \label{sga3:VIIB:2.4.A:eq1}
\end{equation}
is bijective.

The morphism of formal \(k\)-varieties
\begin{equation}
  \theta:G(B)\longrightarrow N(B),
  \qquad
  g\longmapsto g\cdot\sigma\pi(g^{-1})
  \tag{2}
  \label{sga3:VIIB:2.4.A:eq2}
\end{equation}
is then a retraction of \(\sigma'\), and the inverse of \(\mu\) is
\begin{equation}
  g\longmapsto\bigl(\theta(g),\pi(g)\bigr).
  \tag{3}
  \label{sga3:VIIB:2.4.A:eq3}
\end{equation}
If \(p:G\to G/Q\) denotes the quotient projection, then
\[
  p\circ\sigma':N\longrightarrow G/Q
\]
is an isomorphism of formal \(k\)-varieties.

\smallskip
\noindent\emph{Translator's note.} The French re-edition prints
\(\theta\circ\sigma':N\to G/Q\) here. Since
\(\theta:G\to N\), that composite is ill-typed; the isomorphism
with the displayed source and target is the quotient composite
\(p\circ\sigma'\).

\smallskip
In particular, \(N\) is topologically flat over \(k\), by \SGARefLink{sga3:VIIB:1.4.thm}{1.4}
\textup{\SGARefLink{sga3:vii:paragraph:5-1:equation:ii}{(ii)}}. Let
\[
  \alpha:Q\times N\longrightarrow N
\]
and
\[
  \beta:G\times N=N\times Q\times N\longrightarrow N
\]
be the maps defined on points by
\[
  \alpha(q,y)=qyq^{-1},
  \qquad
  \beta(x,q,y)=x\alpha(q,y).
\]
Then the following diagram is commutative:
\begin{center}
  \phantomsection
  \label{sga3:VIIB:2.4.A:eq4-diagram}
  \begin{tikzcd}[row sep=3.2em,column sep=3.5em]
{} \arrow[d,phantom,"\textup{(4)}" description]
  & G\times G
  \arrow[r,"{\operatorname{id}\times\theta}"]
  \arrow[d,"m_G"']
  & G\times N
      \arrow[r,"{\theta\times\operatorname{id}}"]
      \arrow[d,"\beta"]
  & N\times N \arrow[dl,"m_N"] \\
{} & G \arrow[r,"\theta"] & N & {.}
\end{tikzcd}

\end{center}

This can be expressed as follows in terms of the affine algebras
\(A\), \(A_0\), and \(A'\) of \(G\), \(Q\), and \(N\)
\textup{(cf.\ \SGARefLink{sga3:VIIB:5.1.3}{5.1.3} below)}. Let
\[
  \rho':A\longrightarrow A',
  \qquad
  \rho:A\longrightarrow A_0,
  \qquad
  \tau:A_0\longrightarrow A
\]
be the homomorphisms of \(k\)-bialgebras corresponding to
\(\sigma'\), \(\sigma\), and \(\pi\), and put
\[
  I=\operatorname{Ker}(\rho).
\]
By the preceding corollary, \(A'\) is identified with
\[
  A\big/\overline{A\,\tau(J_0)},
\]
where \(J_0\) is the augmentation ideal of \(A_0\).

On the other hand, let \(B\) be the affine algebra of \(G/Q\).
It is the kernel of the pair of morphisms
\phantomsection
\label{sga3:VIIB:2.4.A:affine-equalizer-pair}
\begin{equation*}
  A
  \mathrel{\overset{\tau_1}
    {\underset{(\operatorname{id}_A\widehat\otimes\rho)
      \circ\Delta_A}{\rightrightarrows}}}
  A\widehat\otimes_kA_0,
\end{equation*}
where \(\tau_1(x)=x\widehat\otimes1\). Thus
\[
  B
  =
  \bigl\{
    x\in A
    \mathrel{\big|}
    \Delta_A(x)-x\widehat\otimes1
    \in A\widehat\otimes I
  \bigr\}.
\]
Let
\[
  \gamma=\theta^\natural:A'\longrightarrow A
\]
be the corresponding continuous homomorphism of \(k\)-algebras.
It is a section of \(\rho'\), and it is an isomorphism of profinite
\(k\)-algebras from \(A'\) onto \(B\). Moreover,
\[
  \tau=\pi^\natural
\]
identifies \(A_0\) with a sub-bialgebra of \(A\), namely the affine
algebra of the quotient \(N\backslash G\).

It follows from equations~\eqref{sga3:VIIB:2.4.A:eq1}
and~\eqref{sga3:VIIB:2.4.A:eq3} that there is an
isomorphism of profinite \(k\)-algebras
\begin{equation}
  A'\widehat\otimes_kA_0
  \xrightarrow{\ \sim\ }
  A,
  \qquad
  a'\widehat\otimes a_0
  \longmapsto
  \gamma(a')\tau(a_0),
  \tag{*}
  \label{sga3:VIIB:2.4.A:star}
\end{equation}
whose inverse is the map
\[
  a\longmapsto
  (\rho'\widehat\otimes\rho)\Delta_A(a).
\]

Finally, identify \(A'\) with its image in \(A\) under \(\gamma\),
so that the projection \(A\to A'\) is then \(\gamma\rho'\).
Writing \(\Delta_{A'}\) for the comultiplication of \(A'\), one
deduces from \hyperref[sga3:VIIB:2.4.A:eq4-diagram]{\textup{(4)}}
that
\[
  \Delta_A(A')\subset A\widehat\otimes A'
\]
and that the following diagram is commutative:
\begin{center}
  \phantomsection
  \label{sga3:VIIB:2.4.A:Aprime-comultiplication-square}
  \begin{tikzcd}[row sep=3em,column sep=3.7em]
A' \arrow[r,"\Delta_A"] \arrow[d,equals]
  & A\widehat\otimes A' \arrow[d,"{\gamma\rho'\widehat\otimes
                                      \operatorname{id}}" right] \\
A' \arrow[r,"\Delta_{A'}"]
  & A'\widehat\otimes A'
\end{tikzcd}

\end{center}
Consequently, one also has
\[
  \gamma\rho'\circ c_A=c_{A'},
\]
where \(c_A\) and \(c_{A'}\) are the antipodes of \(A\) and \(A'\),
respectively. On the other hand, writing \(m_A\) for multiplication
in \(A\), equation~\eqref{sga3:VIIB:2.4.A:eq2} gives, for every
\(a\in A\),
\[
  \gamma\rho'(a)
  =
  \bigl(
    m_A
    \circ
    (\operatorname{id}_A\widehat\otimes\tau\rho c_A)
    \circ
    \Delta_A
  \bigr)(a).
\]
